\documentclass[10pt,foldmark,tumble]{leaflet}
\renewcommand*\foldmarkrule{.3mm}
\renewcommand*\foldmarklength{5mm}

\usepackage{amsmath}
\usepackage[T1]{fontenc}
\usepackage{textcomp}
\usepackage{mathptmx}
\usepackage{sans}
\usepackage{pst-barcode}
\usepackage{array}
\usepackage{caption}

\makeatletter
\def\ptmTeX{T\kern-.1667em\lower.5ex\hbox{E}\kern-.075emX\@}
\DeclareRobustCommand{\ptmLaTeX}{L\kern-.3em
        {\setbox0\hbox{T}%
         %\vb@xt@ % :-)
         \vbox to\ht0{\hbox{%
                            \csname S@\f@size\endcsname
                            \fontsize\sf@size\z@
                            \math@fontsfalse\selectfont
                            A}%
                      \vss}%
        }%
        \kern-.12em
        \ptmTeX}
\makeatother
\let\TeX=\ptmTeX
\let\LaTeX=\ptmLaTeX
\usepackage{shortvrb}
\MakeShortVerb{\|}
\usepackage{url}
\usepackage{graphicx}
\usepackage{color}
\definecolor{DESYBlue}{RGB}{0,166,235}
\definecolor{DESYOrange}{RGB}{242,142,0}

%%%%\renewcommand{\descfont}{\normalfont}
\newcommand\Lpack[1]{\textsf{#1}}
\newcommand\Lclass[1]{\textsf{#1}}
\newcommand\Lopt[1]{\texttt{#1}}
\newcommand\Lprog[1]{\textit{#1}}

\newcommand*\defaultmarker{\textsuperscript\textasteriskcentered}

\newcommand{\barcode}[1]{
\begin{pspicture}(1cm,1cm) 
\psbarcode[scalex=0.4,scaley=0.4]{#1}{}{qrcode} 
\end{pspicture}
}

\newcolumntype{C}[1]{>{\centering}m{#1}}
\newcolumntype{L}[1]{>{}m{#1}}

\title{\textcolor{DESYBlue}{\bf Sicherheit@Home}\textcolor{DESYOrange}{\bf.}}

\author{%
\Large \textcolor{DESYBlue}{\bf Deutsches Elektronen-Synchrotron}
}
\date{\textcolor{DESYOrange}{\bf Tag der offenen Tür 2015}}

\CutLine*{1}% Dotted line without scissors
\CutLine{6}%  Dotted line with scissors

%\AddToBackground{1}{%  Background of a small page
%  \put(0,0){\textcolor{Cerulean}{\rule{\paperwidth}{\paperheight}}}}


\AddToBackground{1}{%  Background of a small page
  \put(210,525){\includegraphics[scale=0.1]{logo_desy}}}


\AddToBackground{6}{%  Background of a small page
  \put(0,0){\textcolor{DESYBlue}{\rule{\paperwidth}{\paperheight}}
  }}


%\AddToBackground*{2}{% Background of a large page
%  \put(\LenToUnit{.5\paperwidth},\LenToUnit{.5\paperheight}){%
%    \makebox(0,0)[c]{%
%      \resizebox{.9\paperwidth}{!}{\rotatebox{35.26}{%
%        \textsf{\textbf{\textcolor{DESYOrange}{Tag der offenen Tür 2015}}}}}}}}

\begin{document}

\maketitle

Vorschläge zur Verbesserung Ihrer Sicherheit im Internet

\newpage

\section{Das Internet}

Das Internet ist ein offener Zusammenschluss aus Milliarden von unterschiedlichsten
Teilnehmern. Während der ursprüngliche Zweck des Internets der Austausch von Informationen
unter Wissenschaftlern war, hat heute nahezu jede Organisation eine Internetpräsens und
es tangiert nahezu all unsere Lebensbereiche. Diese Tatsache macht das Internet leider auch
zu einem Werkzeug von unterschiedlichen Angreifern, die versuchen illegal Gewinn aus dieser
Tatsache zu ziehen.
Die Angreifer nutzen unterschiedlichste Taktiken und je nach Ziel des Angreifers kann jeder
Internetteilnehmer von diesen Angriffen betroffen sein:

\begin{itemize}
 \item Server sind häufig mit hoher Bandbreite an das Internet angebunden; das 
 macht diese Systeme als Ziel interessant, da sich darüber schnell Daten verschicken lassen.
 \item Viele Organisationen haben auch interne Computer an das Internet angeschlossen, um 
 zum Beispiel unterschiedliche Filialen zu vernetzen. Diese System enthalten oft geheime Daten, 
 wie Kundenpasswörter oder interne Dokumente.
 \item Private Benutzer haben häufig schwach gesicherte Computer, so dass sie leicht massenhaft
 automatisiert kompromittiert werden können. Die Angreifer können dann Computer weiter ausspionieren, 
 um zum Beispiel Passwörter zu erlangen oder um den Computer für Angriffe auf weitere Systeme 
 zu missbrauchen.
\end{itemize}

\section{Sicherheitsrisiken}

\subsection{Software- und Konfigurationsfehler}

Die hohe Komplexität in heutigen Softwaresystemen und Infrastrukturen führt dazu, dass
Sicherheitslücken entstehen können, die mögliche Einfallstore für Angreifer darstellen. 
Gerade weit verbreitete Software, wie Betriebssysteme, Web- und Datenbankserver werden 
von Angreifern auf Schwachstellen analysiert, die dann den Angriff auf viele Ziele 
ermöglicht, die die gleiche Software einsetzen. Leider haben verschiedene Geheimdienste 
aktiv darauf hingewirkt Sicherheitsmechanismen zu schwächen. Diese Schwachstellen können 
natürlich auch von nicht legitimierten Angreifern ausgenutzt werden.
Konfigurationsfehler, z.B. Standardpasswörter oder veraltete Softwarepakete können 
Angreifern erlauben in interne Netzwerke einzubrechen und von dort weitere Systeme unter
Kontrolle zu bringen.

\subsection{Benutzerfehler}

Der weitaus grössere Anteil der Angriffe nutzt jedoch die Achtlosigkeit der Benutzer aus. 
Bei gezielten Angriffen wird z.B. versucht, sich über Phishing-eMails deren Zugriffsrechte 
anzueignen. Diese eMails werden zum Teil speziell für die Mitarbeiter einer Organisation 
zugeschnitten und gleichen diesen unter Umständen selbst in feinen Details. Um solche eMails 
dennoch entlarven zu können, müssen die Benutzer geschult werden.
Aber auch ungezielt können die Benutzer ungewollt zu Schwachstellen in der Sicherheit
werden: Webbrowser und deren Plugins (z.B. das Flash-Plugin) haben in der Vergangen-
heit häufig Fehler enthalten, die es Angreifern erlaubt, Schadcode auf den Computern der
Anwender zu installieren.
Verlorene oder gestohlene Geräte, z.B. SmartPhones enthalten Passwörter und weitere sen-
sible Daten, die Angreifer leicht auslesen und für weitere Angriffe nutzen können.

\section{Hilfe zur Selbsthilfe}

\begin{figure}
\centering
\includegraphics[width=\textwidth]{HTTPS-Check.jpg/HTTPS-check_003}
\caption{Zertifikat bei HTTPS Verbindungen prüfen}
\label{sslcert}
\end{figure}

Um im Internet nicht schutzlos unterwegs zu sein, sollten Benutzer einige Verhaltensweisen
beherzigen:

\begin{itemize}
 \item Nutzen Sie wenn möglich verschlüsselte Verbindungen (https) zu Webservern, insbesondere
 wenn diese sensible Daten abfragen und vergewissern Sie sich, dass die Zertifikate der Webserver
 gültig sind (siehe Abb. 1).
 \begin{itemize}
  \item[\barcode{bit.ly/1pgk9fU}] HTTPS Everywhere
 \end{itemize}
 
 \item Bei eMails unbekannter Herkunft oder eMails, die Aufforderungen enthalten, Zugangsdaten
 preiszugeben ist besondere Vorsicht geboten. Überprüfen Sie den Absender und eventuell enthaltene
 Links. Öffnen Sie Anhänge nur, wenn die Herkunft der eMail zweifelsfrei geklärt ist!
 
 \item Halten Sie Ihr Betriebssystem und Programme - insbesondere Webbrowser und deren Plugins - 
 immer aktuell!
 
 \item Informieren Sie sich über aktuelle Sicherheitslücken. 
 \begin{itemize}
  \item [\barcode{bit.ly/1Wj6BFs}] BSI für Bürger
 \end{itemize}
 
 \item Nutzen Sie Sicherheitsplugins, die verhindern dass Webseiten unkontrollierbar weitere
 Ressourcen nachladen!
 \begin{itemize}
  \item[\barcode{mzl.la/1Wj85zi}] NoScript Plugin
  \item[\barcode{mzl.la/1XAQMqK}] $\mu{}$Matrix Plugin
 \end{itemize} 

 \item Vermeiden Sie unnötige Preisgabe Ihrer Daten! Schon durch die Verknüpfung Ihrer 
 Verbindungsdaten können detaillierte Profile erstellt werden, die es erlauben Sie gezielt zu 
 manipulieren.
 \begin{itemize}
  \item[\barcode{mzl.la/1nJHRBI}] Tracking kontrollieren mit Ghostery Plugin
  \item[\barcode{mzl.la/1RCMkDS}] Werbenetzwerke erkennen mit Lightbeam Plugin
 \end{itemize}

 \item Verschlüsseln Sie alle sensiblen Daten! Emails, Chatprotokolle und andere Dateien können
 von Angreifern leicht ausgelesen werden, wenn sie nicht verschlüsselt sind. 

 \begin{itemize}
  \item[\barcode{mzl.la/1hrHU3n}] eMail-Verschlüsselung mit Enigmail Plugin für Thunderbird
  \item[\barcode{bit.ly/1E3UvBw}] Sicherheitsvergleich unterschiedlicher Kommunikationsprogramme
 \end{itemize}
  
 \item Geben Sie keine Zugangsdaten über unverschlüsselte WLAN-Netzwerken ein! Mobilgeräte tun
 dies unter Umständen ohne Ihr Zutun, sobald Sie sich verbinden. Benutzen Sie ein VPN, wenn sie 
 öffentliche Hotspots nutzen möchten.
 \begin{itemize}
  \item[\barcode{bit.ly/1PVVVrk}] Erklärung zu VPN
 \end{itemize}

 \item Viele Produkte haben freie Alternativen, die in vielen Fällen eine bessere Sicherheit
 und besseren Datenschutz für die Anwender bieten. Eine Liste von interessanten Programmen finden
 Sie hier:
 \begin{itemize}
 \item[\barcode{bit.ly/1P7djIR}] Freie und sichere Alternativen zu kommerziellen Programmen
 \item[\barcode{bit.ly/1nSJfiZ}] Weitere Empfehlungen zur Sicherheit
 \end{itemize}
\end{itemize}

\vspace{1cm} 
Die Tabelle auf der nächsten Seite stellt eine Auswahl, aus dem Vergleich von \url{https://www.eff.org/secure-messaging-scorecard} dar.

\begin{center}
\tiny{
\rotatebox{90}{
\begin{tabular}{ L{2cm} | C{2cm} | C{2cm} | C{2cm} | C{2cm} | C{2cm} | C{2cm} | C{2cm} }
\label{comparePrograms}

Prog. & Ver\-schlüs\-sel\-te Über\-tra\-gung? & Un\-sicht\-bar beim Pro\-vider? & Iden\-tität be\-kannt? & Frühere Daten sicher, wenn Schlüssel ge\-stohlen? & Sourcen offen? & Programm Design dokumentiert? & Gab es einen zeitnahen code Audit? \tabularnewline
\hline
AIM & \includegraphics[scale=0.1]{pics/haken.png} & \includegraphics[scale=0.1]{pics/nohaken.png} & \includegraphics[scale=0.1]{pics/nohaken.png} & \includegraphics[scale=0.1]{pics/nohaken.png} & \includegraphics[scale=0.1]{pics/nohaken.png} & \includegraphics[scale=0.1]{pics/nohaken.png} & \includegraphics[scale=0.1]{pics/nohaken.png} \tabularnewline[0.5ex]
Facebook chat & \includegraphics[scale=0.1]{pics/haken.png} & \includegraphics[scale=0.1]{pics/nohaken.png} & \includegraphics[scale=0.1]{pics/nohaken.png} & \includegraphics[scale=0.1]{pics/nohaken.png} & \includegraphics[scale=0.1]{pics/nohaken.png} & \includegraphics[scale=0.1]{pics/nohaken.png} & \includegraphics[scale=0.1]{pics/haken.png} \tabularnewline
Skype & \includegraphics[scale=0.1]{pics/haken.png} & \includegraphics[scale=0.1]{pics/nohaken.png} & \includegraphics[scale=0.1]{pics/nohaken.png} & \includegraphics[scale=0.1]{pics/nohaken.png} & \includegraphics[scale=0.1]{pics/nohaken.png} & \includegraphics[scale=0.1]{pics/nohaken.png} & \includegraphics[scale=0.1]{pics/nohaken.png} \tabularnewline
FaceTime & \includegraphics[scale=0.1]{pics/haken.png} & \includegraphics[scale=0.1]{pics/haken.png} & \includegraphics[scale=0.1]{pics/nohaken.png} & \includegraphics[scale=0.1]{pics/haken.png} & \includegraphics[scale=0.1]{pics/nohaken.png} & \includegraphics[scale=0.1]{pics/haken.png} & \includegraphics[scale=0.1]{pics/haken.png} \tabularnewline
WhatsApp & \includegraphics[scale=0.1]{pics/haken.png} & \includegraphics[scale=0.1]{pics/nohaken.png} & \includegraphics[scale=0.1]{pics/nohaken.png} & \includegraphics[scale=0.1]{pics/nohaken.png} & \includegraphics[scale=0.1]{pics/nohaken.png} & \includegraphics[scale=0.1]{pics/nohaken.png} & \includegraphics[scale=0.1]{pics/haken.png} \tabularnewline
Google Hangout & \includegraphics[scale=0.1]{pics/haken.png} & \includegraphics[scale=0.1]{pics/nohaken.png} & \includegraphics[scale=0.1]{pics/nohaken.png} & \includegraphics[scale=0.1]{pics/nohaken.png} & \includegraphics[scale=0.1]{pics/nohaken.png} & \includegraphics[scale=0.1]{pics/nohaken.png} & \includegraphics[scale=0.1]{pics/haken.png} \tabularnewline
iMessage & \includegraphics[scale=0.1]{pics/haken.png} & \includegraphics[scale=0.1]{pics/haken.png} & \includegraphics[scale=0.1]{pics/nohaken.png} & \includegraphics[scale=0.1]{pics/haken.png} & \includegraphics[scale=0.1]{pics/nohaken.png} & \includegraphics[scale=0.1]{pics/haken.png} & \includegraphics[scale=0.1]{pics/haken.png} \tabularnewline
Jitsi+Ostel  & \includegraphics[scale=0.1]{pics/haken.png} & \includegraphics[scale=0.1]{pics/haken.png} & \includegraphics[scale=0.1]{pics/haken.png} & \includegraphics[scale=0.1]{pics/haken.png} & \includegraphics[scale=0.1]{pics/haken.png} & \includegraphics[scale=0.1]{pics/haken.png} & \includegraphics[scale=0.1]{pics/nohaken.png} \tabularnewline
Mailvelope   & \includegraphics[scale=0.1]{pics/haken.png} & \includegraphics[scale=0.1]{pics/haken.png} & \includegraphics[scale=0.1]{pics/haken.png} & \includegraphics[scale=0.1]{pics/nohaken.png} & \includegraphics[scale=0.1]{pics/haken.png} & \includegraphics[scale=0.1]{pics/haken.png} & \includegraphics[scale=0.1]{pics/haken.png} \tabularnewline
Adium+OTR & \includegraphics[scale=0.1]{pics/haken.png} & \includegraphics[scale=0.1]{pics/haken.png} & \includegraphics[scale=0.1]{pics/haken.png} & \includegraphics[scale=0.1]{pics/haken.png} & \includegraphics[scale=0.1]{pics/haken.png} & \includegraphics[scale=0.1]{pics/haken.png} & \includegraphics[scale=0.1]{pics/nohaken.png} \tabularnewline
Pidgin+OTR & \includegraphics[scale=0.1]{pics/haken.png} & \includegraphics[scale=0.1]{pics/haken.png} & \includegraphics[scale=0.1]{pics/haken.png} & \includegraphics[scale=0.1]{pics/haken.png} & \includegraphics[scale=0.1]{pics/haken.png} & \includegraphics[scale=0.1]{pics/haken.png} & \includegraphics[scale=0.1]{pics/haken.png} \tabularnewline
GPGTools & \includegraphics[scale=0.1]{pics/haken.png} & \includegraphics[scale=0.1]{pics/haken.png} & \includegraphics[scale=0.1]{pics/haken.png} & \includegraphics[scale=0.1]{pics/nohaken.png} & \includegraphics[scale=0.1]{pics/haken.png} & \includegraphics[scale=0.1]{pics/haken.png} & \includegraphics[scale=0.1]{pics/nohaken.png} \tabularnewline
Retroshare & \includegraphics[scale=0.1]{pics/haken.png} & \includegraphics[scale=0.1]{pics/haken.png} & \includegraphics[scale=0.1]{pics/haken.png} & \includegraphics[scale=0.1]{pics/haken.png} & \includegraphics[scale=0.1]{pics/haken.png} & \includegraphics[scale=0.1]{pics/haken.png} & \includegraphics[scale=0.1]{pics/nohaken.png} \tabularnewline
Signal / Redphone & \includegraphics[scale=0.1]{pics/haken.png} & \includegraphics[scale=0.1]{pics/haken.png} & \includegraphics[scale=0.1]{pics/haken.png} & \includegraphics[scale=0.1]{pics/haken.png} & \includegraphics[scale=0.1]{pics/haken.png} & \includegraphics[scale=0.1]{pics/haken.png} & \includegraphics[scale=0.1]{pics/haken.png} \tabularnewline
Silent Phone & \includegraphics[scale=0.1]{pics/haken.png} & \includegraphics[scale=0.1]{pics/haken.png} & \includegraphics[scale=0.1]{pics/haken.png} & \includegraphics[scale=0.1]{pics/haken.png} & \includegraphics[scale=0.1]{pics/haken.png} & \includegraphics[scale=0.1]{pics/haken.png} & \includegraphics[scale=0.1]{pics/haken.png} \tabularnewline
Silent Text & \includegraphics[scale=0.1]{pics/haken.png} & \includegraphics[scale=0.1]{pics/haken.png} & \includegraphics[scale=0.1]{pics/haken.png} & \includegraphics[scale=0.1]{pics/haken.png} & \includegraphics[scale=0.1]{pics/haken.png} & \includegraphics[scale=0.1]{pics/haken.png} & \includegraphics[scale=0.1]{pics/haken.png} \tabularnewline
Subrosa & \includegraphics[scale=0.1]{pics/haken.png} & \includegraphics[scale=0.1]{pics/haken.png} & \includegraphics[scale=0.1]{pics/haken.png} & \includegraphics[scale=0.1]{pics/nohaken.png} & \includegraphics[scale=0.1]{pics/haken.png} & \includegraphics[scale=0.1]{pics/haken.png} & \includegraphics[scale=0.1]{pics/haken.png} \tabularnewline
SureSpot & \includegraphics[scale=0.1]{pics/haken.png} & \includegraphics[scale=0.1]{pics/haken.png} & \includegraphics[scale=0.1]{pics/haken.png} & \includegraphics[scale=0.1]{pics/nohaken.png} & \includegraphics[scale=0.1]{pics/haken.png} & \includegraphics[scale=0.1]{pics/haken.png} & \includegraphics[scale=0.1]{pics/nohaken.png} \tabularnewline
Telegram (secret chat) & \includegraphics[scale=0.1]{pics/haken.png} & \includegraphics[scale=0.1]{pics/haken.png} & \includegraphics[scale=0.1]{pics/haken.png} & \includegraphics[scale=0.1]{pics/haken.png} & \includegraphics[scale=0.1]{pics/haken.png} & \includegraphics[scale=0.1]{pics/haken.png} & \includegraphics[scale=0.1]{pics/haken.png} \tabularnewline
TextSecure & \includegraphics[scale=0.1]{pics/haken.png} & \includegraphics[scale=0.1]{pics/haken.png} & \includegraphics[scale=0.1]{pics/haken.png} & \includegraphics[scale=0.1]{pics/haken.png} & \includegraphics[scale=0.1]{pics/haken.png} & \includegraphics[scale=0.1]{pics/haken.png} & \includegraphics[scale=0.1]{pics/haken.png} \tabularnewline
\end{tabular}
} % end rotatebox
} % end tiny
\end{center}
 
\end{document}
